% !TeX root = ../PIFB2.tex
\section{From the Inferential Construction to Physical Questions}\label{sec:conclusion}

The construction now has a single probabilistic center: normalized population laws with selective covariant interactions. The composite Gaussian consensus energy remains an explicit worked specialization. Parent laws, induced interactions and predictive feedback follow by exact operations once their inputs are given. Conditional dynamics produce memory and persistent records; general recursive formation remains open.

The following part asks whether these mathematical structures can carry physical significance. It does not add a physical interpretation to every symbol. Instead it develops the clock and geometry calculations, identifies an exact Lorentz representation, and states what would be needed to make them descriptions of time, space and matter.

\clearpage
\part*{Part II: Physical Structure and Participatory Interpretation}
\addcontentsline{toc}{part}{Part II: Physical Structure and Participatory Interpretation}

